\begin{name}
	{\tenchude}
	{TOÁN 10}
	{LỚP TOÁN THẦY PHÁT}
	{Life always offers you a second chance. It’s called tomorrow}
\end{name}
\Opensolutionfile{ansbook}[ans/ansbookDe5]
\TN
\Opensolutionfile{ans}[ans/ansDe5-TN1]
\begin{ex}%[0D8N1-3]%[Tex hóa đề CK2 - form 2025 - đọt 2 -Huỳnh Thanh Chí]
Từ thành phố $A$ đến thành phố $B$ có 3 con đường, từ thành phố $A$ đến thành phố $C$ có 2 con đường, từ thành phố $B$ đến thành phố $D$ có 2 con đường, từ thành phố $C$ đến thành phố $D$ có 3 con đường, không có con đường nào nối từ thành phố $C$ đến thành phố $B$. Hỏi có bao nhiêu con đường đi từ thành phố $A$ đến thành phố $D$.
\choice
{$6$}
{\True $12$}
{$18$}
{$36$}
\loigiai{Số cách đi từ $A$ đến $D$ bằng cách đi từ $A$ đến $B$ rồi đến $D$ là $3\cdot 2=6$.\\
Số cách đi từ $A$ đến $D$ bằng cách đi từ $A$ đến $C$ rồi đến $D$ là $2\cdot 3=6$.\\
Vậy có $6+6=12$ cách.
}
\end{ex}

\begin{ex}%[0D8N2-1]
Cho $k,n$ là các số nguyên dương, $k \leq n$. Trong các phát biểu sau, phát biểu nào \textbf{sai}?
\choice
{$\mathrm{C}_n^k=\mathrm{C}_n^{n-k}$}
{$\mathrm{C}_n^k=\dfrac{n!}{k!(n-k)!}$}
{$\mathrm{A}_n^k=\dfrac{n!}{(n-k)!}$}
{\True $\mathrm{C}_n^k=\mathrm{A}_n^k \cdot k!$}
\loigiai{
Công thức $\mathrm{C}_n^k=\mathrm{A}_n^k \cdot k!$ sai, công thức đúng phải là $\mathrm{C}_n^k=\dfrac{\mathrm{A}_n^k}{k!}$.
}
\end{ex}

\begin{ex}%[0D8N3-2]
Khai triển $(x+1)^5$.
\choice
{$(x+1)^5=x^5+1$}
{\True $(x+1)^5=x^5+5x^4+10x^3+10x^2+5x+1$}
{$(x+1)^5=x^5+x^4+x^3+x^2+x+1$}
{$(x+1)^5=x^5+4x^4+6x^3+6x^2+4x+1$}
\loigiai{
Ta có $(x+1)^5=x^5+5x^4+10x^3+10x^2+5x+1$.
}
\end{ex}

\begin{ex}  :%[0D6H1-1]
Các nhà toán học cổ đại Trung Quốc đã dùng phân số $\dfrac{22}{7}$ để xấp xỉ số $\pi $. Hãy đánh giá sai số tuyệt đối $\Delta $ của giá trị gần đúng này, biết $3{,}1415<\pi <3{,}1416$.
\choice
{$\Delta <0{,}0012$}
{$\Delta <0{,}0014$}
{\True $\Delta <0{,}0013$}
{$\Delta <0{,}0011$}
\loigiai{
Vì $3{,}1415<\pi<3{,}1416$\\
nên  $\dfrac{22}{7}-3{,}1416<\dfrac{22}{7}-\pi<\dfrac{22}{7}-3{,}1415$.\\
Suy ra $\left|\dfrac{22}{7}-\pi\right|<\dfrac{22}{7}-3{,}1415 \approx 0{,}014$
}
\end{ex}

\begin{ex}%[HK1, THPT Chuyên Hùng Vương - Phú Thọ, 2023-2024]%[Lê Quốc Hiệp, 10-11EX-HK1-2324]%[0D6H4-2]
Mẫu số liệu sau cho biết chiều cao (đơn vị cm) của các bạn trong tổ
\[\begin{array}{llllllll}163 & 159 & 172 & 167 & 165 & 168 & 170 & 161\end{array}\]
Khoảng biến thiên của mẫu số liệu này bằng
\choice
{$11$}
{$10$}
{\True $13$}
{$12$}
\loigiai
{
Sắp xếp mẫu số liệu theo thứ tự không giảm
\[
\begin{array}{llllllll}
159 & 161 & 163 & 165 & 167 & 168 & 170 & 172
\end{array}
\]
Khoảng biến thiên $R=x_8-x_1=172-159=13$.
}
\end{ex}

\begin{ex}%[0D0N1-2]
Tung một đồng xu và một con xúc xắc, nhận được kết quả là mặt xuất hiện trên đồng xu (sấp hay ngửa) và số chấm xuất hiện trên con xúc xắc. Số kết quả có thể xảy ra là
\choice
{$2$}
{$6$}
{$8$}
{\True $12$}
\loigiai{
Đồng xu chỉ có hai mặt sấp hay ngửa.\\
Con xúc xắc số chấm xuất hiện là $1,2,3,4,5,6$.\\
Số kết quả có thể xảy ra là $2\cdot 6=12$.
}
\end{ex}

\begin{ex}%[0D0N1-3]
Gieo ngẫu nhiên $1$ con xúc xắc cân đối, đồng chất $5$ lần. Số phần tử của không gian mẫu là
\choice
{$15625$}
{$11$}
{$30$}
{\True$7776$}
\loigiai{
Số phần tử của không gian mẫu là $n(\Omega)=6^5=7776$.
}
\end{ex}

\begin{ex}%[0H9H3-1]
Vectơ nào dưới đây là một vectơ pháp tuyến của đường thẳng đi qua hai điểm $A(2;3)$ và $B(4;1)$?
\choice
{ $\overrightarrow{n}_{1}=(2;-2)$}
{$\overrightarrow{n}_{2}=(2;-1)$}
{\True$\overrightarrow{n}_{3}=(1;1)$}
{$\overrightarrow{n}_{4}=(1;-2)$}
\loigiai{
Đường thẳng đi qua hai điểm $A(2;3)$ và $B(4;1)$ có vectơ chỉ phương là $\overrightarrow{AB}=(2;-2)$ suy ra một vectơ pháp tuyến là $\overrightarrow{n}=(1;1)$.}
\end{ex}

\begin{ex}%[0H9H3-2]%[Dự án đề kiểm tra Toán Khối 10 CK2 NH23-24-Dot 15- Nguyễn Hoàng Việt]%[THPT Đặng Huy Trứ - Thừa Thiên Huế]
Phương trình tham số của đường thẳng $d$ đi qua điểm $A\,(-1; 2)$ và song song với đường thẳng $\Delta: 3 x-13 y+1=0$ là
\choice
{\True $\heva{&x=-1+13 t \\& y=2+3 t}$}
{$\heva{&x=1+13 t \\& y=-2+3 t}$}
{$\heva{&x=1+3 t \\& y=2-13 t}$}
{$\heva{&x=-1-13 t \\& y=2+3 t}$}
\loigiai{
Từ phương trình đường $\Delta \Rightarrow \overrightarrow{n_{\Delta}}=(3;-13)\Rightarrow \overrightarrow{u_{\Delta}}=(13;3)$.\\
Vì $(d)\parallel \Delta $ nên $\overrightarrow{u_d}=\overrightarrow{u_{\Delta}}=(13;3)$. Lại có đường thẳng $(d)$ đi qua $A\,(-1;2)$ khi đó phương trình tham số của $\Delta $ là
$\heva{&x=-1+13t\\&y=2+3t}$}
\end{ex}

\begin{ex}%[10-11-12EX-HK1-2425]%[Huỳnh Xuân Tín]%[0H9N3-3]
Trong mặt phẳng $Oxy$, cho hai đường thẳng $d_1\colon 2x-y+3=0$ và $d_2\colon x+2y+1=0$. Vị trí tương đối của hai đường thẳng $d_1$ và $d_2$ là
\choice
{\True Vuông góc với nhau}
{Trùng nhau}
{Cắt nhau nhưng không vuông góc với nhau}
{Song song}
\loigiai{Đường thẳng $d_1\colon 2x-y+3=0$ có vectơ pháp tuyến là $\overrightarrow{n}_1=(2;-1)$ và  đường thẳng $d_2\colon x+2y+1=0$ có vectơ pháp tuyến là $\overrightarrow{n}_2=(1;2)$.\\
Ta có  $2\cdot1+(-1)\cdot2=0$ nên $\overrightarrow{n}_1\cdot \overrightarrow{n}_2=0$. Khi đó hai đường thẳng $d_1$ và $d_2$ vuông góc với nhau.}
\end{ex}

\begin{ex}%[0H9N5-1]
Trong mặt phẳng ${Oxy}$, cho elip $ {(E)}:\dfrac{x^{2}}{36}+\dfrac{y^2}{4}=1$. Độ dài trục lớn bằng
\choice
{ \True ${12}$ }
{ ${37}$ }
{ ${36}$ }
{ ${4}$ }
\loigiai{
Ta có $a^2=36 \Rightarrow a=6$. Độ dài trục lớn bằng $2a={12}$.
}\end{ex}

\begin{ex}%[0H9H5-8]
Viết phương trình chính tắc của Parabol biết đường chuẩn có phương trình $x+1=0$.
\choice
{$y^2=2x$}
{\True $y^2=4x$}
{$y=4x^2$}
{$y^2=8x$}
\loigiai{
Phương trình chính tắc của parabol $\left(P\right)\colon y^2=2px$.\\
Đường chuẩn $x+1=0$ suy ra $\dfrac{p}{2}=1$ $\Rightarrow 2p=4$ $\Rightarrow y^2=4x$. Vậy $y^2=4x$.
}
\end{ex}

\TNTF
\begin{ex}%[0D6H4-2]%[CTST - Lớp 10 - Ôn tập cuối học kì 1 - Đề 2]%[Cao Thành Thái]
Cho bảng số liệu sau:\\
\centerline
{
\begin{tabular}{|>{\centering\arraybackslash}p{2cm}|>{\centering\arraybackslash}p{1cm}|>{\centering\arraybackslash}p{1cm}| >{\centering\arraybackslash}p{1cm}| >{\centering\arraybackslash}p{1cm}| >{\centering\arraybackslash}p{1cm}| >{\centering\arraybackslash}p{1cm}|}
\hline
Giá trị & $21$ & $32$ & $18$ & $24$ & $25$ & $26$ \\
\hline
Tần số & $7$ & $6$ & $4$ & $8$ & $6$ & $10$ \\
\hline
\end{tabular}
}
Xét tính đúng sai của các khẳng định dưới đây.
\choiceTF
{Mốt của mẫu số liệu trên là $10$}
{Số trung bình của mẫu số liệu (làm tròn đến hàng phần chục) là $24{,}9$}
{\True Quy tròn giá trị phương sai của mẫu số liệu với độ chính xác $d=0{,}23$ ta được kết quả là $15$}
{\True Sai số tương đối của số quy tròn phương sai đến hàng đơn vị không vượt quá $1\%$}
\loigiai
{
\begin{itemchoice}
\itemch Giá trị có tần số lớn nhất là $26$ nên $M_O = 26$.
\itemch \[\overline{x} = \dfrac{21\cdot 7 + 32\cdot 6 + 18\cdot 4 + 24\cdot 8 + 25\cdot 6 + 26\cdot 10}{40} = \dfrac{199}{8} \approx 24{,}7.\]
\itemch Phương sai của mẫu số liệu là $s^2=15,036\ldots$, quy tròn với độ chính xác $d=0{,}23$ ta được $s^2\approx 15$.\\
\itemch Sai số tương đối của số quy tròn phương sai đến hàng đơn vị là $\dfrac{|15-15,036\ldots|}{15}\approx 0{,}24\% < 1\%$.
\end{itemchoice}
}
\end{ex}

\begin{ex}%[0H9H3-2]
Trong mặt phẳng tọa độ $Oxy$, cho tam giác $ABC$ có $A(1;1)$, $B(2;-2)$, $C(4;2)$.
\choiceTF
{\True Đường thẳng $AB$ vuông góc với đường thẳng $AC$}
{Phương trình đường cao $AH$ là $x+2y+3=0$}
{\True Đường tròn ngoại tiếp tam giác $ABC$ có tâm $I(3;0)$}
{Đường tròn tâm $C$ đi qua điểm $A$ cắt đường thẳng $AB$ tại hai điểm phân biệt}
\loigiai{
\begin{itemchoice}
\itemch Ta có $\overrightarrow{AB}=(1;-3)$, $\overrightarrow{AC}=(3;1)$. Suy ra $\overrightarrow{AB}\cdot\overrightarrow{AC}=0$ nên $AB\perp AC$.
\itemch Vì vectơ $ \overrightarrow{BC}=\left(2;4\right) $ là vectơ pháp tuyến của đường cao $ AH $.\\
Phương trình đường cao $ AH $ là $ x+2y-3=0 $.
\itemch Vì tam giác $ABC$ vuông tại $A$ nên tâm đường tròn ngoại tiếp tam giác $ABC$ là trung điểm $I(3;0)$ của cạnh huyền $BC$.\\
\itemch Đường tròn tâm $C$ đi qua điểm $A$ có bán kính $AC \perp AB$ nên $AB$ là tiếp tuyến của đường tròn đó. Do đó, đường tròn tâm $C$ đi qua điểm $A$ chỉ có một điểm chung với đường thẳng $AB$ là điểm $A$.
\end{itemchoice}
}
\end{ex}

\TNSA
\begin{ex}%[0D8V1-3]
Trong mặt phẳng toạ độ $Oxy$, ở góc phần tư thứ nhất ta lấy $2$ điểm phân biệt; cứ thế ở các góc phần tư thứ hai, thứ ba, thứ tư lần lượt lấy $3,4,5$ điểm phân biệt (các điểm không nằm trên các trục toạ độ). Trong 14 điểm đó ta lấy $2$ điểm bất kỳ và nối chúng lại, hỏi có bao nhiêu đoạn thẳng cắt hai trục toạ độ, biết đoạn thẳng nối 2 điểm bất kì không qua $O$.
\shortans[]{23}
\loigiai{
\begin{center}
\begin{tikzpicture}[>=stealth,line join=round,line cap=round,font=\footnotesize,scale=1]
\draw[->] (-4,0)--(4,0)node[below]{$x$};
\draw[->] (0,-4)--(0,4)node[right]{$y$};
\fill (0,0)node[below left]{$O$}circle(1pt);
\coordinate(A) at (2,3);
\coordinate(B) at (3,2);
\coordinate(C) at (-2,3.5);
\coordinate(D) at (-3,2.5);
\coordinate(E) at (-1,1);
\coordinate(F) at (-3,-1);
\coordinate(G) at (-2,-1.5);
\coordinate(H) at (-2.5,-2.5);
\coordinate(I) at (-3.5,-2);
\coordinate(J) at (2,-1);
\coordinate(K) at (3,-1.5);
\coordinate(M) at (1.5,-2);
\coordinate(L) at (2.5,-3);
\coordinate(N) at (1.5,-3.5);

\fill[black](A)circle (1.5pt);
\fill[black](B)circle (1.5pt);
\fill[black](C)circle (1.5pt);
\fill[black](D)circle (1.5pt);
\fill[black](E)circle (1.5pt);
\fill[black](F)circle (1.5pt);
\fill[black](G)circle (1.5pt);
\fill[black](H)circle (1.5pt);
\fill[black](I)circle (1.5pt);
\fill[black](J)circle (1.5pt);
\fill[black](K)circle (1.5pt);
\fill[black](L)circle (1.5pt);
\fill[black](M)circle (1.5pt);
\fill[black](N)circle (1.5pt);

\draw[line width=0.7pt,black] (B)--(A)--(D)  (G)--(A)--(J);

\end{tikzpicture}
\end{center}
Để chọn $2$ điểm trong $14$ điểm đã cho nối lại cắt hai trục toạ độ thì hai điểm đó phải thuộc hai góc phần tư đối đỉnh với nhau.
\begin{itemize}
\item  Chọn 1 điểm ở góc phần tư thứ I và 1 điểm ở góc phần tư thứ III.\\
Số đoạn thẳng tạo thành là $2\cdot 4=8$.
\item  Chọn 1 điểm ở góc phần tư thứ II và 1 điểm ở góc phần tư thứ IV.\\
Số đoạn thẳng tạo thành là $3 \cdot 5=15$.
\end{itemize}
Theo quy tắc cộng ta có $8+15=23$ đoạn thẳng.}
\end{ex}

\begin{ex}%[0D8V2-5]
Trong kho đèn trang trí đang còn $6$ bóng đèn loại I, $8$ bóng đèn loại II, các bóng đèn đều khác nhau về màu sắc và hình dáng. Lấy ra $5$ bóng đèn bất kỳ. Hỏi có bao nhiêu khả năng xảy ra số bóng đèn loại I nhiều hơn số bóng đèn loại II?
\shortans{$686$}
\loigiai{
Có 3 trường hợp xảy ra:
\begin{itemize}
\item 	TH1: Lấy được $5$ bóng đèn loại I có $\mathrm{C}_6^5$ cách.
\item 	TH2: Lấy được $4$ bóng đèn loại I, $1$ bóng đèn loại II có $\mathrm{C}_6^4 \cdot \mathrm{C}_8^1$ cách.
\item 	TH3: Lấy được $3$ bóng đèn loại I, $2$ bóng đèn loại II có $\mathrm{C}_6^3 \cdot \mathrm{C}_8^2$ cách.
\end{itemize}
Theo quy tắc cộng, ta có: $\mathrm{C}_6^5+\mathrm{C}_6^4 \cdot \mathrm{C}_8^1+\mathrm{C}_6^3 \cdot \mathrm{C}_8^2=686$ cách.
}
\end{ex}

\begin{ex}%[0H9V5-0]
\immini[thm]{Khúc cua của một con đường có dạng hình parabol, điểm đầu vào khúc cua là $A$, điểm cuối là $B$, khoảng cách $AB=400$ m. Đỉnh parabol $(P)$ của khúc của cách đường thẳng $AB$ một khoảng $20$ m và cách đều $A$, $B$ như hình bên. Tìm tham số tiêu của parabol $(P)$ ($1$ đơn vị đo trong mặt phẳng tọa độ tương ứng $1$ m trên thực tế).
}{
\begin{tikzpicture}[line join=round, line cap=round,>=stealth,thick]
\coordinate (A) at (0,-2);
\coordinate (B) at (0,2);
\draw[name path=cong] (A)..controls +(-250:1.7) and +(250:1.7)..(B);
\path[name path=ox] ($(A)!0.5!(B)$)--+(-0.75,0);
\path [name intersections={of=cong and ox,by={C}}];
\draw[dashed](A)--(B)node[pos=0.7, right]{$400$\,m} ($(A)!0.5!(B)$)--(C);
\draw [<-, dashed] (-0.25,0)--(0.5,-0.5) node[right]{$20$\,m};
\draw[fill=black] (A) circle(1.0pt) node[right]{$A$} (B)circle(1.0pt) node[right]{$B$};
\end{tikzpicture}
}

\shortans[]{$1000$}
\loigiai{
Chọn hệ trục tọa độ Oxy sao cho đỉnh của parabol trùng với gốc tọa độ $O(0;0)$ (như hình vẽ).
\begin{center}
\begin{tikzpicture}[line join=round, line cap=round,>=stealth,thick]
\coordinate (A) at (0,-2);
\coordinate (B) at (0,2);
\draw[name path=cong] (A)..controls +(-250:1.7) and +(250:1.7)..(B);
\path[name path=ox] ($(A)!0.5!(B)$)--+(-0.75,0);
\path [name intersections={of=cong and ox,by={C}}];
\draw[dashed](A)--(B)node[pos=0.7, right]{$400$\,m};
\draw [<-, dashed] (-0.25,0)--(0.5,-0.5) node[right]{$20$\,m};
\draw[fill=black] (A) circle(1.0pt) node[right]{$A$} (B)circle(1.0pt) node[right]{$B$} (C) circle(1.0pt) node [below left]{$O$} ($(A)!0.5!(B)$) circle(1.0pt) node [above right]{$H$};
\draw[->](-1,0)--(2,0) node[below]{$x$};
\draw[->]($(C)+(-90:3)$)--($(C)+(90:3)$) node[left]{$y$};
\end{tikzpicture}
\end{center}
Gọi $H$ là hình chiếu của $O$ lên $AB$, khi đó ta chứng minh được $H$ là trung điểm của AB nên $HA=HB=\dfrac{1}{2}AB$.\\
Khoảng cách từ khúc cua đến đường thẳng $AB$ là $OH$.\\
Ta có $AB=400$ m. Suy ra, $HA=HB=400:2=200$ (m); $OH=20$ (m).\\
Do $1$ đơn vị đo trong mặt phẳng tọa độ tương ứng $1$ m trên thực tế thì tọa độ các điểm là $A(20;-200)$ và $B(20;200)$.\\
Gọi phương trình parabol $(P)$ có dạng $y^2=2px$ (với $p>0$).\\
Khi đó $A$, $B$ đều thuộc $(P)$.\\
Thay tọa độ điểm $B$ vào phương trình parabol $(P)$ ta có $200^2=2p\cdot 20\Leftrightarrow p=1\,000$.\\
Vậy parabol $(P)$ có tham số tiêu là $p=1\,000$.
}
\end{ex}

\begin{ex}%[dự án HSA, Huỳnh Xuân Tín]%[0H9V4-2]
Đường tròn $(C)$ có tâm $I$ thuộc đường thẳng $d\colon x+3y+8=0$, đi qua điểm $A\left(-2;1\right)$ và tiếp xúc với đường thẳng $\Delta:3x-4y+10=0$. Tính bán kính đường tròn $(C)$.
\shortans[1]{$5$}
\loigiai{
Ta có $I\in d \Rightarrow I(-3m-8;m)$.
\begin{eqnarray*}
&&AI^2=\left[\mathrm{d}(I,\Delta)\right]^2 \text{ (cùng bằng } R^2) \\
& \Leftrightarrow & \left(-3m-6\right)^2+\left(m-1\right)^2=\dfrac{\left[3(-3m-8)-4m+10\right]^2}{3^2+(-4)^2}\\
& \Leftrightarrow & 10m^2+34m+37=\dfrac{169m^2+364m+196}{25}\\
& \Leftrightarrow & 81m^2+486m+729=0\\
& \Leftrightarrow & m=-3.
\end{eqnarray*}
Do đó, đường tròn $(C)$ có $\heva{&\text{tâm } I(1;-3)\\& \text{bán kính } R=AI=\sqrt{(9-6)^2+(-3-1)^2}=5.}$\\
Vậy phương trình đường tròn là ${\left(x-1\right)}^2+{\left(y+3\right)}^2=25$.
}
\end{ex}

\TL
\begin{ex}
	Tìm số hạng chứa $x^3$ trong khai triển $x(x+\dfrac{2}{x})^4$.
\end{ex}

\begin{ex}%[0H9V4-5]
Trong mặt phẳng với hệ trục tọa độ $Oxy$, cho đường thẳng $\Delta\colon 3x+4y-m=0$. 
\begin{enumerate}
	\item Viết phương trình đường tròn $\left(C\right)$ có tâm $I(-2;1)$ và bán kính $R=5$.
	\item Tìm giá trị dương của tham số $m$ để đường thẳng $\Delta$ cắt đường tròn $(C)$ theo một dây cung có độ dài bằng $6$.
\end{enumerate}
\loigiai{
\begin{center}
\begin{tikzpicture}[declare function={r=2.5;ga=40;},>=stealth,font=\scriptsize,scale=0.8]
%\path (0:0) coordinate (I)
%(ga:r) coordinate (A)
%(130:r) coordinate (B)
\coordinate (I) at (0,0);
\coordinate (A) at (ga:r);
\coordinate (B) at (130:r);
\coordinate (J) at ($(A)!1/2!(B)$);
\draw (I) circle (r)  ;
\draw (I)--(A)--(B)--(I)--(J);
\foreach \x/\g in{A/90,B/90,J/90,I/-90}
\fill[black](\x) circle (2pt)
($(\x)+(\g:5mm)$) node{\small $\x$};
\end{tikzpicture}
\end{center}
Đường tròn $(C)$ có tâm $I(-2;1)$ và bán kính $R=\sqrt{(-2)^2+1^2+20}=5$.\\
Giả sử $\Delta$ cắt $(C)$ tại hai điểm $A$ và $B$. Gọi $J$ là trung điểm của $AB$.\\
Khi đó $\heva{&IJ\perp AB\\&JB=\dfrac{1}{2}AB=3.}$\\
Xét tam giác $IJB$ vuông tại $J$ có $IJ=\sqrt{IB^2-JB^2}=4.$\\
Hay $\mathrm{d}(I; \Delta)=4$. Tương đương
\begin{align*}
&\dfrac{|(-2)\cdot3+4\cdot1-m|}{\sqrt{3^+4^2}}=4\\
\Leftrightarrow&\, |m+2|=4\cdot5=20\\
\Leftrightarrow&\, \hoac{&m+2=20\\&m+2=-20}\Leftrightarrow\hoac{&m=18\\&m=-22\,(\text{không thỏa mãn}).}
\end{align*}
Vậy $m=18$.
}

\end{ex}
\begin{ex}%[Đề phát triển 02]%[Nam Nhất Trần, dự án TeX hóa và phát triển đề tốt nghiệp năm 2025]%[0D0C2-5]
Xung quanh bờ hồ hình tròn, người ta trồng $10$ cây chuối. Người ta chặt ngẫu nhiên $4$ cây chuối.
Tính xác suất để $4$ cây bị chặt không có hai cây nào đứng cạnh nhau.
% \shortans{47}
\loigiai{
\textbf{Cách 1.} \\
Đánh số các cây trên bờ hồ hình tròn từ $1$ đến $10$. \\
Đánh số 4 câu được chọn ngẫu nhiên là $a$, $b$, $c$, $d$ theo thứ tự tăng dần. \\
Số phần tử không gian mẫu là $\mathrm{C}^4_{10}$. \\
Gọi $A$ biến cố thỏa bài toán thì $A$: \lq\lq $1\leq a < b-1 < c-2 < d-3 \leq 7$ \rq\rq.
\begin{itemize}
\item Trường hợp 1. $a=1$. Khi đó $\heva{&b\geq 3 \\&d\leq 9} \Rightarrow \heva{&b-1 \geq 2 \\&d-3 \leq 6.}$ \\
Từ tập $\{2;3;4;5;6\}$, chọn $3$ số và xếp theo tứ tự tăng dần (vào $b-1$, $c-2$, $d-3$), có $\mathrm{C}^3_{5}$ cách. \\
Mỗi cách chọn trên tương ứng với một cách chọn cho $b$, $c$, $d$ thỏa bài toán. \\
Vậy trường hợp này có $\mathrm{C}^3_{5}$ cách.
\item Trường hợp 2. $a\geq 2$. \\
Từ tập $\{2;3;4;5;6;7\}$ chọn 4 số theo và xếp theo thứ tự tăng dần (vào $a$, $b-1$, $c-2$, $d-3$), có $\mathrm{C}^4_{6}$ cách.\\
Mỗi cách chọn trên tương ứng với một cách chọn cho $a$, $b$, $c$, $d$ thỏa bài toán. \\
Vậy trường hợp này có $\mathrm{C}^4_{6}$ cách.
\end{itemize}
Từ hai trường hợp trên suy ra $n(A)=\mathrm{C}^3_{5}+\mathrm{C}^4_{6}=25$ và xác suất của biến cố $A$ là $\mathrm{P}(A)=\dfrac{n(A)}{n(\Omega)}=\dfrac{5}{42}$.
\\
\textbf{Cách 2.} \\
Ta nhắc lại bài toán chia kẹo Euler dạng cơ bản: Phương trình
\[x_1+x_2+...+x_k=n\]
với điều kiện $n\ge k$, có $\mathrm{C}^{k-1}_{n+k-1}$ nghiệm nguyên không âm.
\begin{itemize}
\item Trường hợp 1, giả sử chặt cây $1$. Khi đó cây $2$ và cây $10$ sẽ không bị chặt, vậy ba cây còn lại bị chặt là ba trong số các cây từ $3$ đến $9$. \\
Gọi $x_1$, $x_2$, $x_3$, $x_4$ lần lượt là số lượng cây nằm ở giữa bốn cây bị chặt (chưa bao gồm cây $2$ và cây $10$). \\
Vậy ta có $x_1+x_2+x_3+x_4=4$. Đồng thời $x_1\ge 0$, $x_2\ge 1$, $x_3\ge 1$ và $x_4\ge 0$ (do đã có cây $2$ và cây $10$ đứng cạnh cây $1$ bị chặt, nên $x_1$ và $x_4$ có thể bằng $0$ vẫn thỏa mãn yêu cầu đề bài). \\
Đặt $x_2=y_2+1$, $x_3=y_3+1$, suy ra $x_1+y_2+y_3+x_4=2$ và $x_1,  y_2, y_3, x_4\ge 0$. \\
Áp dụng bài toán chia kẹo Euler, ta có số cách chặt là $\mathrm{C}^3_5=10$ cách.
\item Trường hợp $2$, giả sử cây bị chặt không phải cây $1$. \\
Ta ``trải hình'' thành đường thẳng, duỗi đường tròn tại vị trí cây số $1$. Khi đó, ta xét các cây từ $2$ đến $10$. \\
Ta gọi $x_1$, $x_2$, $x_3$, $x_4$, $x_5$ lần lượt là số cây chưa bị chặt nằm bên trái cây bị chặt số $1$, số cây nằm giữa cây bị chặt $1$ và $2$, số cây nằm giữa cây bị chặt $2$ và $3$, số cây nằm giữa cây bị chặt $3$ và $4$, số cây nằm bên phải cây bị chặt số $4$. \\
Vậy ta có $x_1+x_2+x_3+x_4+x_5=5$ (do từ $2$ đến $10$ có $9$ cây, chặt mất $4$ cây nên còn $5$ cây). \\
Ngoài ra, ta có $x_2, x_3, x_4\ge 1$, đồng thời $x_1\ge 0$ và $x_5\ge 0$. \\
Đặt $x_2=y_2+1$, $x_3=y_3+1$, $x_4=y_4+1$, ta có $x_1+y_2+y_3+y_4+x_5=2$ và $x_1,y_2,y_3,y_4,x_5\ge 0$. \\
Áp dụng bài toán chia kẹo Euler, ta có số cách chặt cây là $\mathrm{C}^4_6=15$ cách.
\end{itemize}
Vậy có $25$ cách chặt thỏa mãn yêu cầu, suy ra xác suất cần tính là $\dfrac{25}{\mathrm{C}^4_{10}}=\dfrac{5}{42}$.
}
\end{ex}
\Closesolutionfile{ans}
\Closesolutionfile{ansbook}
